% =====================================================================
%  E7: heads from the most active decoded words, then an attention filter.
%  Record of the experiment: why it was run, the design decisions, the
%  sets, both checks, and the comparison with Davidson's FV heads.
%  Tables are generated by scripts/e7_tables.py (paper/tables/e7_*.tex).
% =====================================================================
\documentclass[10pt,a4paper]{article}

\usepackage[utf8]{inputenc}
\usepackage[T1]{fontenc}
\usepackage[margin=1.8cm]{geometry}
\usepackage{amsmath,amssymb}
\usepackage{booktabs,tabularx,array}
\usepackage{enumitem}
\usepackage{xcolor}
\usepackage{caption}
\usepackage{float}
\usepackage{graphicx}
\usepackage[colorlinks=true,linkcolor=blue!50!black,urlcolor=blue!60!black]{hyperref}

\setlist{nosep,leftmargin=*}
\setlength{\parindent}{0pt}
\setlength{\parskip}{5pt}
\renewcommand{\arraystretch}{1.12}
\captionsetup{font=small,labelfont=bf}

\title{E7: Heads from the Most Active Decoded Words,\\ Filtered by Attention\\\large Lexical task heads, simple tasks}
\author{Filippo Merlo}
\date{9 November 2026}

\begin{document}
\maketitle
\vspace{-1.5em}

\paragraph{In short.} Heads that decode a task's words but are not among Davidson et al.'s function-vector (FV) heads mostly attend to the first token (BOS); the FV task heads attend to the prompt. E7 turns this into a selection rule: take the heads that decode the task's most active words, then drop those that look mainly at BOS. The filter cuts each set to about a quarter and keeps or raises its effect in check~1, while random subsets of the same size lose almost all of it. The filtered set does about as well as FV sets of similar size, with mostly different heads: it wins on word relations, morphology and translations, FV wins on present--past and the landmark, park and product facts. Removing the words heads also decode on WikiText halves the effect.

% ---------------------------------------------------------------------
\section{Why: what separates FV task heads from other word-decoding heads}

Question: among heads that decode the task's words, what distinguishes the 24 FV task heads (the 30 FV heads of the 14-task run minus the 6 generic ones) from the CC heads that are not FV heads? Per task and head we compared the CC-list match rate, the FV score and its rank in the task, and last-token attention (Table~\ref{tab:profile}). These are whole-study readings, before the held-out split.

\begin{table}[H]\centering\small
% Generated by scripts/e7_tables.py; do not edit.
\begin{tabular}{lrrrr}
\toprule
 & FV task, decode & FV task, no decode & CC-only & all, L13--31 \\
\midrule
heads per task & 5 & 19 & 17 & 548 \\
layer & 15 & 15 & 19 & 22 \\
attention on BOS & 0.25 & 0.39 & 0.80 & 0.81 \\
attention on task words & 0.21 & 0.05 & 0.02 & 0.01 \\
attention on the input & 0.15 & 0.09 & 0.04 & 0.04 \\
output alike across tasks & 0.29 & 0.38 & 0.52 & 0.54 \\
CC-list match rate & 0.81 & 0.00 & 0.90 & 0.00 \\
FV score rank in task & 20 & 50 & 404 & 537 \\
\bottomrule
\end{tabular}

\caption{Medians per (task, head). FV task heads split by whether they decode the task's CC list (rate $\geq$0.3); CC-only = in the task's CC top 20 by match rate, not an FV head. Attention at the last prompt token, averaged over tasks. Source: \texttt{results/davidson\_fv\_diff/fv\_vs\_cc.json}.}
\label{tab:profile}
\end{table}

\begin{itemize}
\item \textbf{Decoding does not predict the FV score.} Over the 574 decoding (task, head) pairs of layers 13--31 (rate $\geq$0.4), the match rate barely follows the FV rank (Spearman 0.09). CC-only heads decode the words more often than FV heads, yet rank around 400 in the task.
\item \textbf{Where the head looks does.} Attention on BOS is the strongest single predictor ($\rho = -0.49$), then attention on the rest of the template, task words and the input (0.46, 0.38, 0.35). Among decoding pairs, 44 of the 116 with BOS attention below 0.4 are in the task's FV top 30, against 20 of the other 458.
\item CC-only heads spend about 80\% of their attention on BOS, like an average head of these layers; they sit later (layer 19 vs 15) and their output varies less across tasks. FV task heads that decode the words attend to the task words and the input.
\item \textbf{Davidson's set is one set for all 14 tasks}, ranked by the mean score. 15 CC-only heads rank in some task's FV top 30 (L16H18 first on English--French) but on 1--3 tasks only; FV heads are in the top 30 on a median of 5. Narrow heads never enter the FV set.
\end{itemize}

% ---------------------------------------------------------------------
\section{Design}

All choices were fixed before patching. E4's protocol, J-head only, rule style, the 14 simple tasks.

\begin{enumerate}
\item \textbf{Word lists} (per task, on the choosing half; E4's word rules: a head's top 10 at the last prompt token, decoded, lowercased, alphabetic, $\geq$3 letters, no stop words, Zipf $\geq$3). A word's \emph{activity} = the share of (prompt, head) pairs whose top 10 holds it, over the 500 pooled prompts of the five rule wordings.
  \begin{itemize}
  \item \emph{Raw}: the task's 10 most active words. No cross-task contrast, no separate common words.
  \item \emph{Clean}: the same after removing the 50 words heads decode most on 1{,}000 WikiText-103 passages (the cut of CC-pooled-clean).
  \end{itemize}
\item \textbf{Heads}: the CC-pooled rule, a list token in the head's top 10 on $\geq$10\% of the 500 pooled prompts.
\item \textbf{Attention filter} ($+$F): drop the heads whose mean last-token attention on BOS over the task's rule prompts of the choosing half is above 0.5, i.e.\ heads giving more than half of their attention to BOS rather than the prompt.
\item \textbf{Controls}, 5 draws each, patched with the task means in check~1 at the $\alpha$ of the filtered set: \emph{random-drop}, as many heads as the filtered set drawn from the unfiltered set (does the filter pick the right heads, or does a smaller set just do better?); \emph{layer-matched random}, the filtered set's count per layer drawn outside it (as E4).
\item \textbf{Checks} (E4): check~1, items answered correctly with one shot, patched in their zero-shot prompts; check~2, items answered wrongly with one shot, patched in their 1-shot prompts. Rule means of the choosing half and the control mean, ten $\alpha$ ($-10, -4, 1, 2, 3, 4, 5, 6, 8, 10$), replies $\leq$16 tokens, flexible score. Accuracy at the cross-fitted $\alpha$ (the best one on the other fold), mean of the two folds. Every answer saved per item. E4's CC, CC-20/40 and FV-20/40/$N_t$ were run on the same items and folds.
\end{enumerate}

\paragraph{Decisions taken along the way.}
\begin{itemize}
\item \emph{Raw list risk.} Without a contrast, the raw lists could have come out alike on every task (\emph{what, which, https}). They were printed and checked before any GPU work; they did not (Sec.~\ref{sec:sets}). Both lists were kept.
\item \emph{Pooled vs per-wording rule.} Pooling over 500 prompts lets a head pass on one wording alone if it is active on about half of that wording's prompts; ``$\geq$10\% in at least one wording'' would be looser. The pooled rule was kept for comparability with CC; heads passing on one wording only are counted (Table~\ref{tab:sizes}).
\item \emph{Threshold}: 0.5, with set sizes at 0.4 and 0.6 reported but not used.
\end{itemize}

% ---------------------------------------------------------------------
\section{Word lists and sets}\label{sec:sets}

\begin{table}[H]\centering\scriptsize
% Generated by scripts/e7_tables.py; do not edit.
\begin{tabular}{>{\raggedright\arraybackslash}p{2.3cm}>{\raggedright\arraybackslash}p{5.9cm}>{\raggedright\arraybackslash}p{5.9cm}r}
\toprule
Task & raw list (fold 1) & clean list (fold 1) & fold 2 same \\
\midrule
antonym & \emph{opposite, what, which, word, instead, etc, because, http, opposing, words} & \emph{opposite, word, opposing, words, therefore, definition, reverse, synonyms, whats, name} & 9/9 \\
synonym & \emph{what, which, etc, synonyms, amongst, https, word, maybe, whilst, their} & \emph{synonyms, word, words, similar, synonym, definition, name, whats, english, syn} & 9/10 \\
singular--plural & \emph{because, what, single, which, https, http, name, noun, etc, multiple} & \emph{single, name, noun, multiple, plural, either, linguistic, therefore, names, named} & 10/10 \\
present--past & \emph{what, etc, english, https, verb, http, which, grammar, name, maybe} & \emph{english, verb, grammar, name, linguistic, question, verbs, called, word, how} & 9/9 \\
country--capital & \emph{city, cities, which, country, what, town, countries, answer, mph, international} & \emph{city, cities, country, town, countries, answer, international, geography, downtown, capital} & 10/10 \\
country--currency & \emph{currency, currencies, country, which, monetary, money, what, countries, dollar, international} & \emph{currency, currencies, country, monetary, money, countries, dollar, international, financial, answer} & 10/10 \\
landmark--country & \emph{country, which, countries, england, canada, germany, what, france, location, australia} & \emph{country, countries, england, canada, germany, france, location, australia, city, international} & 9/9 \\
park--country & \emph{country, which, countries, what, england, canada, location, australia, park, located} & \emph{country, countries, england, canada, location, australia, park, located, france, national} & 9/9 \\
product--company & \emph{company, companies, which, manufacturer, what, product, http, etc, maybe, microsoft} & \emph{company, companies, manufacturer, product, microsoft, manufacturers, products, named, name, owned} & 10/7 \\
next-item & \emph{next, which, what, january, july, september, november, october, february, mph} & \emph{next, january, july, september, november, october, february, sequence, question, items} & 4/4 \\
prev-item & \emph{previous, what, preceding, prior, answer, question, item, which, next, because} & \emph{previous, preceding, prior, answer, question, item, next, items, sequence, january} & 9/9 \\
English--French & \emph{french, france, english, paris, what, which, translation, http, etc, their} & \emph{french, france, english, paris, translation, italian, european, translations, language, linguistic} & 9/8 \\
English--German & \emph{german, germany, germans, english, which, http, what, berlin, european, french} & \emph{german, germany, germans, english, berlin, european, french, translation, linguistic, language} & 9/10 \\
English--Spanish & \emph{spanish, english, what, mexican, http, hispanic, translation, bilingual, spain, which} & \emph{spanish, english, mexican, hispanic, translation, bilingual, spain, language, mexico, translations} & 9/10 \\
\bottomrule
\end{tabular}

\caption{Lists of fold 1; last column: words also in fold 2's lists (raw/clean). Removed WikiText words include \emph{which, what, http, https, mph, maybe, etc, instead, because}.}
\label{tab:lists}
\end{table}

The raw lists are roughly half task words and half generic words: \emph{what} and \emph{which} are on all 14, \emph{http} on 7. That is close to CC-pooled's own mix (5 task + 5 common words). After the WikiText cut no word is on more than 6 lists. The folds agree on 9--10 of the 10 words except next-item (different months come up).

\begin{table}[H]\centering\small
\resizebox{\linewidth}{!}{% Generated by scripts/e7_tables.py; do not edit.
\begin{tabular}{lrrrrrrrrrr}
\toprule
Task & raw & +filter & at 0.4 & at 0.6 & BOS med. & one wording & clean & +filter & BOS med. & one wording \\
\midrule
antonym & 124--125 & 42--43 & 28--30 & 56 & 0.67 & 26 & 84--90 & 30--31 & 0.62 & 12--14 \\
synonym & 133--134 & 33--35 & 16--19 & 52 & 0.68 & 30--34 & 81--82 & 23 & 0.66 & 20--21 \\
singular--plural & 118--120 & 38--39 & 22--23 & 50--51 & 0.64 & 23--24 & 84--85 & 38--39 & 0.57 & 11--12 \\
present--past & 106--111 & 26--30 & 11--13 & 41--43 & 0.67 & 26--28 & 58--60 & 23 & 0.64 & 8--10 \\
country--capital & 116--118 & 32 & 16--17 & 39--41 & 0.74 & 21--23 & 92--95 & 26--27 & 0.71 & 15--17 \\
country--currency & 105--106 & 26--28 & 14--16 & 36--39 & 0.76 & 15 & 80--81 & 22--24 & 0.73 & 6 \\
landmark--country & 103--107 & 31--34 & 20--24 & 44--48 & 0.67 & 4--9 & 76--87 & 26--32 & 0.61 & 1--5 \\
park--country & 108--125 & 34--39 & 24--27 & 41--49 & 0.71 & 14--19 & 84--99 & 30--34 & 0.64 & 9--14 \\
product--company & 130--135 & 33--37 & 21--23 & 42--46 & 0.74 & 16--19 & 75--87 & 29--32 & 0.62 & 3--9 \\
next-item & 82--130 & 22--28 & 10--13 & 36--46 & 0.68 & 18--23 & 76--101 & 24--27 & 0.65 & 9--10 \\
prev-item & 101--111 & 35--37 & 26--29 & 45--46 & 0.67 & 3--8 & 90--101 & 32--35 & 0.66 & 5--6 \\
English--French & 104--121 & 28--29 & 20--21 & 36--39 & 0.70 & 14--18 & 64--66 & 20--21 & 0.63 & 9--10 \\
English--German & 72--88 & 17--19 & 12--14 & 25--28 & 0.73 & 11 & 62 & 22 & 0.63 & 4--5 \\
English--Spanish & 92--98 & 29--30 & 19--21 & 41--42 & 0.66 & 15--20 & 61--67 & 21--22 & 0.64 & 9--10 \\
\bottomrule
\end{tabular}
}
\caption{Heads per set (range over the two folds); median BOS attention of the unfiltered set; heads reaching 10\% on only one of the five wordings.}
\label{tab:sizes}
\end{table}

The filter keeps a third or less of every set: most word-decoding heads look mainly at BOS. The filtered sets (17--43 heads) are about the size of E4's CC-20/40 and FV-20/40. About one raw head in five passes on a single wording, fewer for the clean list. The raw sets share about 85\% of their heads with E4's CC sets, the clean ones about 70\%.

% ---------------------------------------------------------------------
\section{Results}

\subsection{Check 1}

\begin{table}[H]\centering\small
\resizebox{\linewidth}{!}{% Generated by scripts/e7_tables.py; do not edit.
\begin{tabular}{lrrrrrrrrrr}
\toprule
Task & items & unpatched & raw & raw+F & clean & clean+F & drop raw & layer raw & drop clean & layer clean \\
\midrule
antonym & 100/100 & 0.00 & \textbf{0.62} & 0.46 & 0.21 & 0.17 & 0.05 & 0.00 & 0.08 & 0.00 \\
synonym & 100/100 & 0.03 & \textbf{0.37} & 0.30 & 0.09 & 0.12 & 0.09 & 0.03 & 0.06 & 0.04 \\
singular--plural & 91/88 & 0.00 & 0.79 & \textbf{0.88} & 0.54 & 0.56 & 0.12 & 0.00 & 0.38 & 0.01 \\
present--past & 100/100 & 0.00 & 0.18 & 0.21 & 0.21 & \textbf{0.38} & 0.01 & 0.00 & 0.02 & 0.00 \\
country--capital & 84/86 & 0.00 & 0.76 & \textbf{0.87} & 0.60 & 0.56 & 0.41 & 0.00 & 0.11 & 0.01 \\
country--currency & 67/65 & 0.00 & \textbf{0.32} & 0.32 & 0.12 & 0.09 & 0.01 & 0.00 & 0.06 & 0.00 \\
landmark--country & 100/100 & 0.01 & 0.02 & 0.05 & 0.00 & 0.03 & 0.13 & 0.03 & \textbf{0.20} & 0.01 \\
park--country & 100/100 & 0.00 & 0.33 & \textbf{0.42} & 0.14 & 0.16 & 0.16 & 0.01 & 0.12 & 0.03 \\
product--company & 100/100 & 0.07 & 0.33 & \textbf{0.46} & 0.22 & 0.33 & 0.14 & 0.08 & 0.36 & 0.13 \\
next-item & 58/63 & 0.00 & 0.16 & \textbf{0.17} & 0.01 & 0.05 & 0.07 & 0.01 & 0.02 & 0.01 \\
prev-item & 53/50 & 0.00 & 0.20 & \textbf{0.32} & 0.05 & 0.10 & 0.02 & 0.00 & 0.02 & 0.00 \\
English--French & 100/100 & 0.00 & 0.25 & \textbf{0.40} & 0.02 & 0.07 & 0.00 & 0.00 & 0.02 & 0.01 \\
English--German & 100/100 & 0.00 & 0.41 & \textbf{0.54} & 0.16 & 0.20 & 0.00 & 0.00 & 0.01 & 0.00 \\
English--Spanish & 100/100 & 0.00 & \textbf{0.44} & 0.12 & 0.11 & 0.01 & 0.04 & 0.01 & 0.00 & 0.00 \\
\midrule
\textbf{mean} &  & 0.01 & 0.37 & \textbf{0.40} & 0.18 & 0.20 & 0.09 & 0.01 & 0.10 & 0.02 \\
control mean &  &  & 0.04 & 0.03 & 0.01 & 0.01 & -- & -- & -- & -- \\
\bottomrule
\end{tabular}
}
\caption{Check 1, task accuracy at the cross-fitted $\alpha$, mean of the two folds. $+$F = BOS filter; drop / layer = random-drop and layer-matched random (mean of 5 draws, at the $\alpha$ of the filtered set). Items: patched items in fold 1/2. Control mean: accuracy with the control task's mean.}
\label{tab:check1}
\end{table}

\begin{itemize}
\item \textbf{The filter keeps the heads that carry the task.} Both lists do as well or better after losing three quarters of their heads, better on 10 of 14 tasks for each. Random subsets of the same size drawn from the same unfiltered sets keep almost nothing; layer-matched random sets nothing.
\item The filter hurts antonym and English--Spanish.
\item \textbf{The WikiText cut halves the effect.} Clean is below raw on nearly every task; only present--past gains. As with CC-pooled-clean, the generic words bring heads that work.
\item Control means stay low: the gains are task-specific.
\end{itemize}

\begin{table}[H]\centering\small
\resizebox{\linewidth}{!}{% Generated by scripts/e7_tables.py; do not edit.
\begin{tabular}{lrrrrrrrrr}
\toprule
Task & items & unpatched & raw+F & CC & CC-20 & CC-40 & FV-20 & FV-40 & FV-$N_t$ \\
\midrule
antonym & 100/100 & 0.00 & 0.46 & 0.42 & 0.01 & 0.09 & 0.40 & 0.38 & \textbf{0.49} \\
synonym & 100/100 & 0.03 & 0.30 & 0.34 & 0.10 & 0.19 & 0.13 & 0.28 & \textbf{0.55} \\
singular--plural & 91/88 & 0.00 & \textbf{0.88} & 0.69 & 0.25 & 0.27 & 0.42 & 0.62 & 0.85 \\
present--past & 100/100 & 0.00 & 0.21 & 0.02 & 0.03 & 0.04 & \textbf{0.87} & 0.68 & 0.85 \\
country--capital & 84/86 & 0.00 & 0.87 & 0.75 & 0.77 & 0.74 & 0.88 & 0.92 & \textbf{0.94} \\
country--currency & 67/65 & 0.00 & 0.32 & 0.29 & 0.14 & 0.14 & 0.14 & 0.11 & \textbf{0.47} \\
landmark--country & 100/100 & 0.01 & 0.05 & 0.92 & 0.69 & 0.95 & 0.94 & \textbf{0.95} & 0.93 \\
park--country & 100/100 & 0.00 & 0.42 & 0.73 & 0.51 & 0.60 & 0.87 & \textbf{0.89} & 0.88 \\
product--company & 100/100 & 0.07 & 0.46 & 0.34 & 0.51 & 0.40 & 0.66 & 0.77 & \textbf{0.91} \\
next-item & 58/63 & 0.00 & 0.17 & 0.26 & 0.03 & 0.16 & 0.02 & 0.26 & \textbf{0.70} \\
prev-item & 53/50 & 0.00 & \textbf{0.32} & 0.26 & 0.03 & 0.07 & 0.04 & 0.03 & 0.19 \\
English--French & 100/100 & 0.00 & \textbf{0.40} & 0.40 & 0.00 & 0.01 & 0.00 & 0.03 & 0.21 \\
English--German & 100/100 & 0.00 & \textbf{0.54} & 0.39 & 0.00 & 0.01 & 0.00 & 0.00 & 0.01 \\
English--Spanish & 100/100 & 0.00 & 0.12 & \textbf{0.38} & 0.00 & 0.00 & 0.00 & 0.05 & 0.25 \\
\midrule
\textbf{mean} &  & 0.01 & 0.40 & 0.44 & 0.22 & 0.26 & 0.38 & 0.43 & \textbf{0.59} \\
control mean &  &  & 0.03 & 0.06 & 0.07 & 0.09 & 0.02 & 0.02 & 0.04 \\
\bottomrule
\end{tabular}
}
\caption{Check 1, raw $+$F against E4's sets on the same items: CC (CC-pooled, $\sim$120 heads), CC and FV ranked sets of 20 and 40 heads, and FV-$N_t$ (as many FV-ranked heads as CC).}
\label{tab:refs}
\end{table}

\begin{itemize}
\item \textbf{Against FV of similar size}, raw $+$F (32 heads on average) sits between FV-20 and FV-40, and above them on 9 and 8 tasks. It wins on antonym, synonym, singular--plural, country--currency, prev-item and the translations; FV wins on present--past, the landmark, park and product facts and next-item.
\item \textbf{Against CC}, with a quarter of its heads, raw $+$F is ahead on 9 tasks and behind on 5. CC's higher mean comes almost entirely from landmark--country; without it the two are level (0.42 vs 0.41).
\item Same-size CC sets (CC-20/40, ranked by match rate) are far below: the filter is a better way to shrink a decoding set than ranking by decoding.
\item FV-$N_t$ remains the best set overall.
\item \textbf{Landmark--country collapses on every E7 set.} The patched model answers with a format, not a country (\emph{Country: Germany}, \emph{Country of origin: Italy}, \emph{A Journey Through Time}). The raw set shares 89 of CC's 130 heads; the 41 others came from CC's words \emph{maybe, https, etc} and seem to prevent this. The same collapse appeared with the E3 group-contrast lists.
\end{itemize}

\subsection{Check 2}

\begin{table}[H]\centering\small
\resizebox{\linewidth}{!}{% Generated by scripts/e7_tables.py; do not edit.
\begin{tabular}{lrrrrrrrrrr}
\toprule
Task & items & unpatched & raw & raw+F & clean & clean+F & CC & FV-20 & FV-40 & FV-$N_t$ \\
\midrule
antonym & 100/100 & 0.00 & 0.30 & 0.34 & 0.34 & 0.35 & 0.35 & 0.30 & \textbf{0.37} & 0.35 \\
synonym & 100/100 & 0.00 & 0.16 & 0.16 & 0.09 & 0.14 & 0.17 & 0.11 & 0.17 & \textbf{0.18} \\
singular--plural & 11/15 & 0.00 & \textbf{0.79} & 0.77 & 0.76 & 0.74 & 0.77 & 0.64 & 0.76 & 0.65 \\
present--past & 30/31 & 0.00 & 0.63 & 0.78 & 0.57 & 0.70 & 0.64 & \textbf{0.89} & 0.77 & 0.82 \\
country--capital & 14/13 & 0.00 & 0.77 & 0.74 & 0.80 & \textbf{0.81} & 0.77 & 0.67 & 0.75 & 0.74 \\
country--currency & 31/34 & 0.00 & 0.21 & 0.12 & 0.11 & 0.15 & \textbf{0.23} & 0.17 & 0.16 & 0.17 \\
landmark--country & 100/87 & 0.00 & 0.17 & 0.19 & 0.18 & 0.21 & 0.32 & \textbf{0.33} & 0.31 & 0.32 \\
park--country & 22/23 & 0.00 & 0.18 & 0.26 & 0.16 & 0.20 & 0.33 & 0.39 & \textbf{0.41} & 0.38 \\
product--company & 84/75 & 0.00 & 0.19 & 0.21 & 0.13 & 0.21 & 0.16 & 0.20 & 0.22 & \textbf{0.22} \\
next-item & 54/50 & 0.00 & 0.22 & 0.33 & 0.20 & 0.22 & 0.26 & 0.26 & 0.46 & \textbf{0.61} \\
prev-item & 59/63 & 0.00 & 0.06 & 0.15 & 0.07 & 0.10 & 0.09 & 0.10 & 0.29 & \textbf{0.32} \\
English--French & 100/100 & 0.00 & 0.20 & 0.22 & 0.19 & 0.20 & 0.16 & 0.21 & 0.29 & \textbf{0.36} \\
English--German & 100/100 & 0.00 & 0.22 & \textbf{0.25} & 0.24 & 0.25 & 0.21 & 0.16 & 0.24 & 0.17 \\
English--Spanish & 100/100 & 0.00 & 0.21 & 0.25 & \textbf{0.26} & 0.24 & 0.24 & 0.19 & 0.25 & 0.22 \\
\midrule
\textbf{mean} &  & 0.00 & 0.31 & 0.34 & 0.29 & 0.32 & 0.33 & 0.33 & 0.39 & \textbf{0.39} \\
control mean &  &  & 0.10 & 0.14 & 0.08 & 0.10 & 0.13 & 0.08 & 0.12 & 0.12 \\
\bottomrule
\end{tabular}
}
\caption{Check 2 (1-shot prompts answered wrongly), task accuracy at the cross-fitted $\alpha$. Random controls not run in this check.}
\label{tab:check2}
\end{table}

As in E4 the sets are close. The filter again helps slightly for both lists (10 and 12 of 14 tasks), mostly on the sequences, present--past and park--country. Two caveats: on prev-item the control mean beats the task mean for raw $+$F (0.49 vs 0.15), so that gain is not task-specific; on singular--plural the control mean is high for every set, since the 1-shot example already carries the task. Several tasks have few items in this check (11--15 on singular--plural and country--capital, 22--34 on park, currency and present--past), so single-task differences are coarse.

% ---------------------------------------------------------------------
\section{The filtered heads and the FV heads}

\begin{table}[H]\centering\small
\resizebox{\linewidth}{!}{% Generated by scripts/e7_tables.py; do not edit.
\begin{tabular}{lrrrrrrrrr}
\toprule
Task & raw+F & FV-20: shared & chance & no decode & filtered & FV-40: shared & chance & no decode & filtered \\
\midrule
antonym & 42.5 & 5.5 & 2.7 & 13.0 & 1.5 & 16.5 & 5.2 & 22.0 & 1.5 \\
synonym & 34.0 & 6.5 & 2.0 & 12.5 & 1.0 & 12.0 & 3.8 & 24.5 & 3.5 \\
singular--plural & 38.5 & 8.5 & 2.5 & 11.5 & 0.0 & 13.0 & 4.3 & 26.5 & 0.5 \\
present--past & 28.0 & 4.5 & 1.0 & 15.5 & 0.0 & 6.5 & 2.1 & 32.0 & 1.5 \\
country--capital & 32.0 & 9.0 & 2.5 & 9.0 & 2.0 & 12.5 & 4.9 & 24.5 & 3.0 \\
country--currency & 27.0 & 2.5 & 0.9 & 17.5 & 0.0 & 4.5 & 2.2 & 34.0 & 1.5 \\
landmark--country & 32.5 & 4.0 & 1.4 & 16.0 & 0.0 & 7.5 & 2.9 & 31.5 & 1.0 \\
park--country & 36.5 & 6.0 & 2.3 & 14.0 & 0.0 & 8.5 & 3.6 & 29.0 & 2.5 \\
product--company & 35.0 & 6.0 & 2.0 & 14.0 & 0.0 & 8.5 & 4.2 & 31.5 & 0.0 \\
next-item & 25.0 & 3.0 & 1.2 & 16.5 & 0.5 & 5.0 & 2.8 & 34.0 & 1.0 \\
prev-item & 36.0 & 4.0 & 1.9 & 15.5 & 0.5 & 8.0 & 3.7 & 31.5 & 0.5 \\
English--French & 28.5 & 3.0 & 0.9 & 16.0 & 1.0 & 4.5 & 1.8 & 33.0 & 2.5 \\
English--German & 18.0 & 2.0 & 0.7 & 18.0 & 0.0 & 2.5 & 1.3 & 37.0 & 0.5 \\
English--Spanish & 29.5 & 1.5 & 1.0 & 17.5 & 1.0 & 3.5 & 2.2 & 35.0 & 1.5 \\
\midrule
\textbf{mean} & 31.6 & 4.7 & 1.7 & 14.7 & 0.5 & 8.0 & 3.2 & 30.4 & 1.5 \\
\bottomrule
\end{tabular}
}
\caption{Heads of raw $+$F and of FV-20 / FV-40, per task, mean of the two folds. Shared: in both. Chance: expected shared heads if the FV set were drawn at random within its layers. No decode: FV heads not in the unfiltered raw set. Filtered: FV heads in the raw set dropped by the BOS filter.}
\label{tab:overlap}
\end{table}

\begin{itemize}
\item \textbf{Mostly different heads, from the same layers} (median layer 15 for FV, 14 for raw $+$F). Overlap is 2--3 times chance and largest where both work (capitals, singular--plural, synonym); near chance on the translations, where only raw $+$F carries the task.
\item \textbf{The filter and the FV score agree; decoding is what differs.} About three quarters of the FV heads do not decode any of the task's 10 most active words. Of those that do, the filter removes very few: 69\% of FV-20 heads attend to BOS at most half of the time.
\end{itemize}

\paragraph{Which threshold would keep the decoding FV heads?} For each task, the largest BOS attention among the FV heads that are in the raw set (worst fold), and the raw $+$F set size at that threshold (Table~\ref{tab:threshold}). Keeping \emph{all} FV heads, decoding or not, would need a threshold up to 0.96, i.e.\ no filter.

\begin{table}[H]\centering\small
% Generated by scripts/e7_tables.py; do not edit.
\begin{tabular}{lrrrr}
\toprule
Task & FV-20: threshold & set size & FV-40: threshold & set size \\
\midrule
antonym & 0.53 & 47--48 & 0.53 & 47--48 \\
synonym & 0.52 & 35--37 & 0.87 & 54--96 \\
singular--plural & 0.39 & 19--21 & 0.54 & 22--41 \\
present--past & 0.47 & 8--20 & 0.58 & 37--42 \\
country--capital & 0.62 & 37--43 & 0.73 & 43--56 \\
country--currency & 0.47 & 4--22 & 0.87 & 30--69 \\
landmark--country & 0.47 & 27--33 & 0.66 & 50--55 \\
park--country & 0.50 & 27--39 & 0.79 & 51--80 \\
product--company & 0.43 & 25--29 & 0.47 & 25--34 \\
next-item & 0.53 & 26--27 & 0.58 & 27--42 \\
prev-item & 0.52 & 15--40 & 0.52 & 16--40 \\
English--French & 0.90 & 10--84 & 0.90 & 35--84 \\
English--German & 0.38 & 4--12 & 0.54 & 4--25 \\
English--Spanish & 0.80 & 17--67 & 0.92 & 28--87 \\
\bottomrule
\end{tabular}

\caption{BOS threshold that keeps every decoding FV head (max over the two folds) and the raw $+$F set size at it (the two folds). The filter used 0.5.}
\label{tab:threshold}
\end{table}

About 0.53 keeps every decoding FV-20 head on 11 of 14 tasks (median 0.47). The rest come from a few heads that decode the task's words but still look mostly at BOS, mostly in layers 17--24 (L17H8 and L20H6 on French, L19H27 and L24H2 on Spanish, L18H9 on country--currency, L10H5 on synonym, L21H14 on capitals). FV-40's decoding heads spread higher (median 0.56); keeping them all would bring back nearly the whole raw set.

% ---------------------------------------------------------------------
\section{Takeaways}

\begin{enumerate}
\item Among heads that decode a task's words, \textbf{the ones that look at the prompt carry the task}; the ones that look at BOS add little. This holds on held-out items, against random subsets of the same size.
\item A small decoding set chosen this way \textbf{matches causally ranked FV sets of similar size} without any causal step, but on different tasks and with mostly different heads: the two selections are complementary.
\item The generic words are needed: the WikiText cut removes heads that work.
\end{enumerate}

\paragraph{Caveats.} One run per setting, greedy, no intervals. Attention is averaged over each task's prompts per head; decoding and the FV score are per task. The motivation analysis (Sec.~1) was run on whole-study readings, before the split. Landmark--country fails for a format reason the flexible score does not accept. Patched items are fewer than 100 on some tasks, and as few as 11--15 in check~2 (Tables~\ref{tab:check1}, \ref{tab:check2}).

\paragraph{Open next steps.}
\begin{itemize}
\item Union of raw $+$F and FV-20 ($\sim$50 heads): if the two sets are complementary, it should beat both where each one wins.
\item Raw $+$F at a 0.55 threshold (keeps nearly all decoding FV-20 heads). Chosen after seeing FV, so a secondary result.
\item FV sets of exactly the filtered set's size per task and fold, for a strict size match.
\item What the non-decoding FV heads read under J-head (present--past, landmark--country).
\end{itemize}

\paragraph{Reproduce.}
\begin{itemize}
\item \texttt{scripts/e3/fv\_vs\_cc\_profile.py}: the head profile of Sec.~1.
\item \texttt{scripts/e7\_sets.py --choose A|B}: lists, sets and controls of each fold.
\item \texttt{scripts/e7\_patch.sh}: patching (GPU, about 75 min on 4 A100).
\item \texttt{scripts/e7\_report.py}: accuracies, \texttt{results/e7\_attfilter/report.md}.
\item \texttt{scripts/e7\_tables.py}: these tables and \texttt{results/e7\_attfilter/fv\_compare.json}.
\end{itemize}

\end{document}
